% rel_framework_thermo.tex
% Agent 3: Relativistic Thermodynamics and Causality Framework (BDNK)
% Conventions: (-,+,+,+) signature, c kept explicit, see RELATIVISTIC_CONVENTIONS.md
% Causality: BDNK first-order formalism, see BDNK_CONVENTIONS.md

\section{Relativistic Thermodynamics and Causality Framework}
\label{sec:rel-thermo-framework}

The classical theory of hydrodynamic and hydromagnetic stability rests on
thermodynamic identities and transport equations that are valid only in
the non-relativistic regime.  When fluid velocities approach $c$, when
thermal energy becomes comparable to rest-mass energy, or when
gravitational fields are strong, every one of these building blocks must
be replaced by its relativistic counterpart.  This section establishes
the relativistic thermodynamic framework and, critically, identifies the
causality constraints that any physically admissible equation of state or
transport law must satisfy.

We adopt the BDNK (Bemfica--Disconzi--Noronha--Kovtun) formalism for
dissipative relativistic hydrodynamics.  Unlike the Israel--Stewart (IS)
theory, which promotes dissipative fluxes to independent dynamical
degrees of freedom with relaxation equations, BDNK achieves causal
propagation at \emph{first order} in the gradient expansion through a
judicious choice of hydrodynamic frame.  The result is a theory with
\emph{only} the conservation laws as evolution equations, no spurious
relaxation modes, and a rigorous proof of strong hyperbolicity.

%------------------------------------------------------------------
\subsection{Relativistic Thermodynamics}
\label{subsec:rel-thermo}

We work in the local rest frame of a fluid element and denote by
$\varepsilon$ the total energy density (including rest-mass energy),
by $p$ the isotropic pressure, by $n$ the baryon number density, by
$s$ the entropy density, and by $T$ the temperature.

\subsubsection{First Law}
The relativistic first law of thermodynamics, expressed per unit
volume, reads
\begin{equation}
  \mathrm{d}\varepsilon
  = T\,\mathrm{d} s
    + \frac{\varepsilon + p}{n}\,\mathrm{d} n \,.
  \label{eq:rel-first-law}
\end{equation}
Here $(\varepsilon+p)/n$ is the relativistic chemical potential
(enthalpy per particle); it reduces to the rest-mass energy $m c^{2}$
plus small thermal corrections in the non-relativistic limit.

\subsubsection{Enthalpy Density and Specific Enthalpy}
The \emph{relativistic enthalpy density} is defined as
\begin{equation}
  w \equiv \varepsilon + p \,.
  \label{eq:enthalpy-density}
\end{equation}
This quantity replaces $\rho c^{2}$ as the effective inertia of the
fluid in relativistic dynamics: the perfect-fluid energy-momentum
tensor is
\begin{equation}
  T^{\mu\nu} = \frac{w}{c^{2}}\,u^{\mu}\,u^{\nu} + p\,g^{\mu\nu}\,,
  \label{eq:perfect-Tmunu}
\end{equation}
so that $w/c^{2}$ plays the role of an effective mass density.

The \emph{relativistic specific enthalpy} (dimensionless) is
\begin{equation}
  h \equiv \frac{\varepsilon + p}{\rho_{0}\,c^{2}}
         = \frac{w}{\rho_{0}\,c^{2}}\,,
  \label{eq:specific-enthalpy}
\end{equation}
where $\rho_{0} = m\,n$ is the rest-mass density.  In the
non-relativistic limit $h \to 1 + (\hat\varepsilon + p)/(\rho_{0}c^{2})$,
where $\hat\varepsilon$ is the internal (thermal) energy density; the
departure from unity measures the relativistic correction.

\subsubsection{Gibbs--Duhem Relation}
Thermodynamic consistency requires the Gibbs--Duhem relation
\begin{equation}
  \mathrm{d} p = s\,\mathrm{d} T + n\,\mathrm{d}\mu \,,
  \label{eq:gibbs-duhem}
\end{equation}
where $\mu = (\varepsilon + p)/n - T s/n$ is the chemical potential per
baryon.  This identity is central to deriving relativistic buoyancy
criteria (Section~\ref{subsec:rel-buoyancy}).

\subsubsection{Equation of State}
A complete equation of state (EOS) specifies $p$ as a function of two
independent thermodynamic variables.  Common choices are:
\begin{itemize}
  \item $p = p(\varepsilon, n)$ — convenient for numerical relativity
        and baryon-conserving flows;
  \item $p = p(\varepsilon, s)$ — natural for isentropic (adiabatic) flows;
  \item $p = p(\rho_0, T)$ — useful for comparison with the classical
        Chandrasekhar framework.
\end{itemize}
The relativistic Euler equations close once the EOS is specified.  The
additional requirement of \emph{causality} constrains the functional
form, as we discuss next.


%------------------------------------------------------------------
\subsection{Lorentz Invariance of Thermodynamic Entropy}
\label{subsec:entropy-lorentz}

The entire relativistic thermodynamic framework employed in this book
rests on the Lorentz invariance of the thermodynamic entropy~$S$.
While this is widely assumed in the literature,
Gavassino~\citep{GavassinoLorentz2021} provided a rigorous proof from
first principles, requiring only two assumptions: (i)~there exists at
least one inertial frame in which the second law $\dot{S} \ge 0$
holds, and (ii)~the microscopic dynamics is governed by a Lorentz-invariant
field-theory Lagrangian.  The proof proceeds by showing that
\emph{adiabatic accelerations} (infinitely slow, mechanically driven
velocity changes) leave the rest-mass invariant to leading order in
the coupling strength, so that the entropy change
$\delta S \sim \epsilon^{2} \ge 0$ for both signs of the coupling
constant~$\epsilon$, which forces $\delta S = 0$.  By van Kampen's
argument~\citep{vanKampen1968}, $S = S(M)$ depends only on the
Lorentz-invariant rest mass, establishing the Lorentz scalar nature
of~$S$.

This result guarantees that the covariant entropy current
$s^{\mu} = s\,n\,u^{\mu} + \ldots$ used throughout this book
transforms correctly under Lorentz boosts, and that the second law
$\nabla_{\mu}s^{\mu} \ge 0$ is a frame-independent statement.

%------------------------------------------------------------------
\subsection{Causality Constraints on the Equation of State}
\label{subsec:eos-causality}

\subsubsection{Sound Speed Bound}
Small adiabatic perturbations of a perfect fluid propagate at the
\emph{adiabatic sound speed}
\begin{equation}
  \cs^{2} = \left(\frac{\partial p}{\partial\varepsilon}\right)_{\!s}\,.
  \label{eq:sound-speed}
\end{equation}
Causality demands that no signal travel faster than light:
\begin{equation}
  \boxed{\;0 \;\le\; \cs^{2} \;\le\; c^{2}\;}\,.
  \label{eq:sound-causality}
\end{equation}
The lower bound $\cs^{2} \ge 0$ is thermodynamic stability (Le
Chatelier's principle): a medium with $\cs^{2} < 0$ is mechanically
unstable to collapse.  The upper bound $\cs^{2} \le c^{2}$ is the
relativistic causality constraint; it is automatically satisfied for
all known forms of matter but must be checked for phenomenological EOS
models.

\subsubsection{Le Chatelier's Principle}
More precisely, thermodynamic stability requires
\begin{equation}
  \left(\frac{\partial p}{\partial\varepsilon}\right)_{\!s} > 0
  \qquad\text{and}\qquad
  \left(\frac{\partial p}{\partial\varepsilon}\right)_{\!n} > 0 \,.
  \label{eq:le-chatelier}
\end{equation}
The first condition ensures mechanical stability (real sound speed);
the second ensures that an increase in energy at constant composition
raises the pressure.  Violation of either signals a thermodynamic
instability that is unrelated to the hydrodynamic instabilities studied
in the Chandrasekhar framework but can trigger gravitational collapse
in compact objects.

\subsubsection{Ideal Relativistic Gas}
For an ideal gas with adiabatic index $\Gamma$,
\begin{equation}
  \cs^{2}
  = \frac{\Gamma\,p}{\varepsilon + p}
  = \frac{\Gamma\,p}{w}\,.
  \label{eq:cs-ideal}
\end{equation}
Two important limits:
\begin{enumerate}
  \item \textbf{Ultrarelativistic limit} ($\varepsilon \gg \rho_{0}c^{2}$,
    radiation-dominated).  Here $p = \varepsilon/3$ and $\Gamma = 4/3$,
    giving
    \begin{equation}
      \cs^{2} \;\longrightarrow\; \frac{c^{2}}{3}
      \qquad\text{(ultrarelativistic)}\,.
      \label{eq:cs-ultrarel}
    \end{equation}
    This is the well-known result for a photon gas or
    ultrarelativistic plasma.

  \item \textbf{Non-relativistic limit} ($k_{B}T \ll mc^{2}$).
    With $\varepsilon \approx \rho_{0}c^{2}$ and $p \approx n k_{B}T$,
    \begin{equation}
      \cs^{2} \;\longrightarrow\;
      \frac{\gamma\,k_{B}\,T}{m} \;\ll\; c^{2}
      \qquad\text{(non-relativistic)}\,,
      \label{eq:cs-nonrel}
    \end{equation}
    where $\gamma$ is the classical ratio of specific heats.  This
    recovers the standard result used throughout the Chandrasekhar
    treatment.
\end{enumerate}


%------------------------------------------------------------------
\subsection{The Zeroth Law and Tolman--Ehrenfest Equilibrium}
\label{subsec:zeroth-law}

Before discussing transport laws, we must first establish what
\emph{thermal equilibrium} means in a relativistic setting.  This is
a subtler question than in Newtonian physics, because a moving body's
temperature is not uniquely defined---the long-standing ``Planck--Ott
controversy'' concerns whether a moving body appears hotter or cooler.

\citet{GavassinoZerothLaw2020} has resolved this controversy by
revisiting the zeroth law of thermodynamics within special relativity.
The key results, essential for the foundations of this book, are:

\subsubsection{Ambiguity of temperature in relativity}
For an isolated system with energy $E$ and additional constants of
motion $X_A$ (e.g.\ conserved charges, momentum), the entropy
$S = S(E, X_A)$ defines the temperature as
$T^{-1} = (\partial S/\partial E)|_{X_A}$.  However, any other set
of conserved quantities $Y_B(E, X_A)$ would yield a different
temperature $T'^{-1} = (\partial S/\partial E)|_{Y_B}$.  This
ambiguity is resolved by the zeroth law: a ``wall permeable only to
heat'' (WPOH) conserves all Noether charges $Q_A$ arising from
internal symmetries, thereby selecting the correct variables and
removing the ambiguity.

\subsubsection{Covariant equation of state}
For an isolated relativistic body with total four-momentum $p^\mu$
and Pauli--Lubanski pseudovector $W^\mu$ (the spin), the covariant
equation of state takes the form
\begin{equation}
  S = S(M, Q_A, W^2)\,,\qquad
  M = \sqrt{-p^\mu p_\mu}\,,
  \label{eq:cov-eos}
\end{equation}
where $M$ is the invariant mass.  The temperature is then
\begin{equation}
  \frac{1}{T} = \frac{\partial S}{\partial M}\bigg|_{Q_A, W^2}\,,
  \label{eq:rel-temperature}
\end{equation}
which is a Lorentz scalar---it is defined in the rest frame of the
body and does not transform under boosts.

\subsubsection{Resolution of the Planck--Ott controversy}
Two bodies are in thermal equilibrium if and only if they:
\begin{enumerate}
  \item[(a)] are at rest with respect to each other, and
  \item[(b)] have the same rest-frame temperature $T$.
\end{enumerate}
Any configuration in which two interacting bodies maintain a relative
velocity for infinite time is a \emph{non-ergodic} metastable state,
for which equilibrium thermodynamics does not apply.  The apparent
``temperatures'' measured by Planck-type or Ott-type thermometers in
relative motion correspond to different conserved quantities being held
fixed during the measurement, and do not represent true thermal equilibrium.

\subsubsection{The Tolman--Ehrenfest condition}
In a stationary gravitational field with timelike Killing vector
$\xi^\mu$, the equilibrium condition becomes
\begin{equation}
  T\,\sqrt{-\xi_\mu\xi^\mu} = \text{const}\,,
  \label{eq:tolman-ehrenfest}
\end{equation}
which is the Tolman--Ehrenfest law: hotter regions are at lower
gravitational potential.  This condition is fundamental for:
\begin{itemize}
  \item Defining the adverse temperature gradient (relative to the
    Tolman--Ehrenfest baseline) in the relativistic B\'enard problem
    (Chapter~II);
  \item Understanding the equilibrium temperature profile in rotating
    neutron stars (Chapter~III);
  \item Correctly formulating the Schwarzschild criterion for
    convective instability in general-relativistic stellar interiors
    (Chapter~VI).
\end{itemize}

The derivation of~\eqref{eq:tolman-ehrenfest} from the maximum entropy
principle is straightforward: in equilibrium, the four-vector
$\beta^\mu = u^\mu/T$ must be a Killing vector (ensuring
$\nabla_\mu s^\mu = 0$).  Since the flow is also along $\xi^\mu$,
this implies $u^\mu/T \propto \xi^\mu$, giving~\eqref{eq:tolman-ehrenfest}.

%------------------------------------------------------------------
\subsection{Thermal Conduction: From Fourier to BDNK}
\label{subsec:heat-conduction}

The treatment of heat conduction illustrates, in the most transparent
way, the necessity of causal transport theory in the relativistic
regime and the elegance of the BDNK resolution.

\subsubsection{Classical Fourier Law (Acausal)}
The Fourier law
\begin{equation}
  \vect{q} = -\kappa\,\nabla T
  \label{eq:fourier}
\end{equation}
combined with the energy equation leads to the diffusion equation
$\partial T/\partial t = \chi\,\nabla^{2}T$ (with
$\chi = \kappa/c_{v}$ the thermal diffusivity).  This equation is
\emph{parabolic}: a temperature perturbation propagates
instantaneously to all points in space.  In the non-relativistic
limit this is an excellent approximation; in the relativistic regime
it is \emph{acausal} and therefore inadmissible.

The relativistic generalization of Fourier's law in the Landau frame
(Eckart theory) is likewise acausal and, moreover, suffers from
generic short-wavelength instabilities that render it unusable even
for numerical computation.

\subsubsection{BDNK Resolution: Causality Through Frame Choice}
\label{subsubsec:BDNK-heat}

The BDNK formalism resolves the acausality of Fourier's law
\emph{without} introducing relaxation equations.  The key insight is
that the Eckart and Landau frame choices artificially constrain
the energy-momentum tensor decomposition, removing degrees of
freedom that are needed to make the first-order system
well-posed.

In the BDNK general frame, the heat flux 4-vector $Q^{\mu}$
(satisfying $Q^{\mu}u_{\mu} = 0$) is given by
\begin{equation}
  Q^{\mu}
  = -\kappa\,T\,\Delta^{\mu\alpha}\,\partial_{\alpha}\!\left(\frac{\mu}{T}\right)
    + \mathcal{F}^{\mu}[\varepsilon_{1}, \beta_{1}, \alpha_{1}]\,,
  \label{eq:BDNK-heat}
\end{equation}
where $\Delta^{\mu\nu} = g^{\mu\nu} + u^{\mu}u^{\nu}/c^{2}$ is the
spatial projection tensor, $\kappa \ge 0$ is the thermal conductivity,
and $\mathcal{F}^{\mu}$ collects the \emph{frame correction terms}
that depend on the BDNK frame coefficients $\varepsilon_{1}$,
$\beta_{1}$, and $\alpha_{1}$.  These frame coefficients modify
the off-equilibrium definitions of energy density and pressure:
\begin{equation}
  E = \varepsilon + \varepsilon_{1}\,,\qquad
  P = p + \Pi_{\text{bulk}}\,,
  \label{eq:BDNK-EP}
\end{equation}
where $\varepsilon_{1}$ is a first-order correction
\begin{equation}
  \varepsilon_{1}
  = -\zeta_{1}\,\nabla_{\mu}u^{\mu}
    - \beta_{1}\,\frac{u^{\mu}\partial_{\mu}T}{T}
    - \alpha_{1}\,u^{\mu}\partial_{\mu}\!\left(\frac{\mu}{T}\right).
  \label{eq:eps1-BDNK}
\end{equation}
In the Landau or Eckart frames, specific choices force $Q^{\mu} = 0$
or $\varepsilon_{1} = 0$, thereby removing the degrees of freedom
needed for a well-posed initial-value problem.  In the BDNK general
frame, these coefficients are free to be chosen such that the
\emph{full system} of conservation equations,
\begin{equation}
  \nabla_{\mu}T^{\mu\nu} = 0\,,\qquad
  \nabla_{\mu}(n\,u^{\mu}) = 0\,,
  \label{eq:BDNK-conservation}
\end{equation}
forms a \emph{strongly hyperbolic} system.  No additional evolution
equations for $Q^{\mu}$ are needed: the conservation laws alone are
the complete equations of motion.

\subsubsection{Comparison with the Telegraph Equation}
In the Israel--Stewart approach, the Cattaneo-type relaxation leads
to the telegraph equation
$\tau_{q}\,\partial^{2}T/\partial t^{2} + \partial T/\partial t
= \chi\,\nabla^{2}T$, which achieves causality by converting the
parabolic diffusion equation into a hyperbolic wave equation with a
finite propagation speed $v_{T} = \sqrt{\chi/\tau_{q}}$.

The BDNK approach achieves the same physical goal by a fundamentally
different mechanism.  The linearized system for temperature
perturbations takes the form of a \emph{first-order} dissipative
equation
\begin{equation}
  \frac{\partial T}{\partial t}
  = \kappa_{\text{eff}}\,\nabla^{2}T\,,
  \label{eq:BDNK-thermal-lin}
\end{equation}
where the effective thermal diffusivity $\kappa_{\text{eff}}$
incorporates the BDNK frame coefficients.  Although this appears
parabolic, the coupling to the momentum and continuity equations
through the frame terms renders the \emph{full coupled system}
strongly hyperbolic.  The characteristic speeds of this system are
determined by the roots of the characteristic polynomial of the
first-order conservation equations, and these roots are guaranteed
to be real and bounded by $c$ when the BDNK conditions are
satisfied.  There is no separate ``thermal wave speed''
$v_{T} = \sqrt{\chi/\tau_{q}}$; instead, the maximum signal speed
emerges from the coupled hydrodynamic system.

\subsubsection{Entropy Production}
The second law of thermodynamics requires non-negative entropy
production.  In the BDNK framework, the entropy current is
\begin{equation}
  S^{\mu} = s\,u^{\mu} - \frac{\mu}{T}\,J^{\mu}
            + \frac{Q^{\mu}}{T}\,,
  \label{eq:BDNK-entropy-current}
\end{equation}
where $J^{\mu} = n\,u^{\mu}$ is the baryon current.  The entropy
production rate satisfies
\begin{equation}
  \nabla_{\mu}S^{\mu}
  = \frac{\Pi^{\mu\nu}\Pi_{\mu\nu}}{2\,\eta\,T}
    + \frac{\Pi_{\text{bulk}}^{2}}{\zeta\,T}
    + \frac{Q^{\mu}Q_{\mu}}{\kappa\,T^{2}}
  \;\ge\; 0
  \label{eq:BDNK-entropy-production}
\end{equation}
provided $\eta \ge 0$, $\zeta \ge 0$, and $\kappa \ge 0$.  The BDNK
frame coefficients $(\varepsilon_{1}, \beta_{1}, \alpha_{1})$ do not
enter the entropy production at leading order---they contribute only
at second order in gradients and thus do not constrain the sign of
$\nabla_{\mu}S^{\mu}$.


%------------------------------------------------------------------
\subsection{Viscosity: BDNK First-Order Treatment}
\label{subsec:rel-viscosity}

The dissipative part of the energy-momentum tensor includes both
\emph{bulk viscous pressure} $\Pi_{\text{bulk}}$ and the
\emph{shear stress tensor} $\Pi^{\mu\nu}$.  In the BDNK formalism,
these are given by first-order constitutive relations with no
relaxation equations required.

\subsubsection{Shear Viscosity}
The shear stress tensor $\Pi^{\mu\nu}$ (symmetric, traceless,
transverse: $\Pi^{\mu\nu}u_{\nu} = 0$, $\Pi^{\mu}{}_{\mu} = 0$) is
given directly by the first-order gradient expansion:
\begin{equation}
  \boxed{\;
  \Pi^{\mu\nu} = -2\,\eta\,\sigma^{\mu\nu}
  \;}
  \label{eq:BDNK-shear}
\end{equation}
where $\eta \ge 0$ is the (dynamic) shear viscosity coefficient and
$\sigma^{\mu\nu}$ is the shear tensor:
\begin{equation}
  \sigma^{\mu\nu}
  = \Delta^{\mu\alpha}\Delta^{\nu\beta}
    \!\left(
      \frac{1}{2}(\nabla_{\alpha}u_{\beta} + \nabla_{\beta}u_{\alpha})
      - \frac{1}{3}\,\Delta_{\alpha\beta}\,\nabla_{\gamma}u^{\gamma}
    \right).
  \label{eq:thermo-shear-tensor}
\end{equation}
This is formally identical to the Navier--Stokes constitutive
relation.  The crucial difference from the acausal Eckart/Landau
first-order theories is that in the BDNK framework, the fluid
4-velocity $u^{\mu}$ is defined in a \emph{general frame} (neither
Landau nor Eckart), and the frame coefficients in the energy and
pressure ensure that the coupled system is hyperbolic.

No shear relaxation time $\tau_{\pi}$ is needed.  The Israel--Stewart
equation $\tau_{\pi}\,u^{\alpha}\nabla_{\alpha}\pi^{\langle\mu\nu\rangle}
+ \pi^{\mu\nu} = 2\eta\,\sigma^{\mu\nu}$ is replaced entirely by the
algebraic relation~\eqref{eq:BDNK-shear} together with the BDNK frame
conditions.

\subsubsection{Bulk Viscosity}
The bulk viscous pressure is
\begin{equation}
  \Pi_{\text{bulk}} = -\zeta\,\nabla_{\mu}u^{\mu}\,,
  \label{eq:BDNK-bulk}
\end{equation}
where $\zeta \ge 0$ is the bulk viscosity coefficient and
$\theta \equiv \nabla_{\mu}u^{\mu}$ is the expansion scalar.  Again,
no bulk relaxation time $\tau_{\Pi}$ is required; causality is
ensured by the BDNK frame coefficients.

\subsubsection{Full BDNK Energy-Momentum Tensor}
Collecting the pieces, the full dissipative energy-momentum tensor in
the BDNK general frame is
\begin{equation}
  T^{\mu\nu}
  = \frac{E}{c^{2}}\,u^{\mu}u^{\nu}
    + P\,\Delta^{\mu\nu}
    + \frac{2\,Q^{(\mu}\,u^{\nu)}}{c}
    + \Pi^{\mu\nu}\,,
  \label{eq:BDNK-Tmunu}
\end{equation}
where $E = \varepsilon + \varepsilon_{1}$ and
$P = p + \Pi_{\text{bulk}}$ include the first-order corrections.  The
evolution of the system is governed \emph{solely} by the conservation
laws~\eqref{eq:BDNK-conservation}; there are no additional relaxation
equations for $Q^{\mu}$, $\Pi_{\text{bulk}}$, or $\Pi^{\mu\nu}$.

\subsubsection{Entropy Production from Viscosity}
Entropy production from shear viscosity:
\begin{equation}
  \sigma_{s}^{(\pi)} = \frac{\Pi^{\mu\nu}\Pi_{\mu\nu}}{2\,\eta\,T}
  \;\ge\; 0\,.
  \label{eq:entropy-shear}
\end{equation}
Entropy production from bulk viscosity:
\begin{equation}
  \sigma_{s}^{(\Pi)} = \frac{\Pi_{\text{bulk}}^{2}}{\zeta\,T}
  = \frac{\zeta\,(\nabla_{\mu}u^{\mu})^{2}}{T} \;\ge\; 0\,.
  \label{eq:entropy-bulk}
\end{equation}


%------------------------------------------------------------------
\subsection{BDNK Causality Conditions}
\label{subsec:BDNK-causality}

The central achievement of the BDNK formalism is that causality and
well-posedness are established through conditions on the transport
and frame coefficients, without introducing new dynamical degrees
of freedom.

\subsubsection{Strong Hyperbolicity}
The system of conservation equations~\eqref{eq:BDNK-conservation}
with the BDNK constitutive relations is \emph{strongly hyperbolic}
(and hence the initial-value problem is well-posed) if and only if
the following conditions hold simultaneously:

\begin{enumerate}
  \item \textbf{Positive shear viscosity:}
  \begin{equation}
    \eta > 0\,.
    \label{eq:BDNK-eta-cond}
  \end{equation}

  \item \textbf{Positive bulk--shear combination:}
  \begin{equation}
    \zeta + \tfrac{2}{3}\,\eta > 0\,.
    \label{eq:BDNK-bulk-shear-cond}
  \end{equation}

  \item \textbf{Frame coefficient inequalities:} The BDNK frame
  coefficients $(\varepsilon_{1}, \beta_{1}, \alpha_{1})$ must
  satisfy a set of coupled nonlinear inequalities that ensure all
  characteristic speeds of the system remain real and bounded by $c$.
  Schematically, for a single-component fluid with thermal
  conductivity $\kappa$, the conditions take the form
  \begin{equation}
    \mathcal{C}(\eta,\,\zeta,\,\kappa;\,\varepsilon_{1},\,\beta_{1},\,\alpha_{1};\,\varepsilon,\,p,\,w,\,\cs^{2}) > 0\,,
    \label{eq:BDNK-coupled-ineq}
  \end{equation}
  where $\mathcal{C}$ denotes the set of constraints derived from the
  requirement that all eigenvalues of the principal symbol of the
  system be real and distinct (strong hyperbolicity).
\end{enumerate}

\subsubsection{Characteristic Speeds}
\label{subsubsec:BDNK-char-speeds}
The characteristic speeds of the BDNK system are determined by the
roots of the characteristic polynomial of the first-order
conservation equations.  For a single-component relativistic fluid,
this is a polynomial whose degree equals the number of independent
hydrodynamic variables (e.g., $\varepsilon$, $n$, $u^{i}$).

In contrast to Israel--Stewart, which introduces extra dynamical
variables ($q^{\mu}$, $\pi^{\mu\nu}$, $\Pi$) and thereby increases
the degree of the characteristic polynomial (producing additional
``relaxation modes''), BDNK has \emph{only} the physical
hydrodynamic modes:
\begin{itemize}
  \item \textbf{Sound modes}: speeds modified by viscosity and thermal
    conductivity through the transport and frame coefficients, but
    bounded by $c$ when the BDNK conditions hold.
  \item \textbf{Shear (diffusive) modes}: at low wavenumber $k$ the
    dispersion relation is $\omega = -i\,\eta\,k^{2}/(w/c^{2}) + O(k^{3})$;
    at high $k$ the modes become propagating with speeds bounded by $c$.
  \item \textbf{Thermal mode}: similarly diffusive at low $k$ and causal
    at high $k$, with no independent ``thermal wave speed.''
\end{itemize}
The absence of spurious modes is a significant advantage for stability
analysis: every mode of the BDNK system has a direct physical
interpretation, and there are no unphysical growing modes that can
contaminate numerical simulations.

\subsubsection{Comparison with Israel--Stewart Causality Conditions}
The IS theory requires three strictly positive relaxation times:
\begin{equation}
  \tau_{q} > 0\,,\qquad
  \tau_{\Pi} > 0\,,\qquad
  \tau_{\pi} > 0\,,
  \tag{IS causality}
  \label{eq:IS-causality-old}
\end{equation}
together with signal-speed bounds $v_{T} \le c$, $v_{\Pi} \le c$,
$v_{\pi} \le c$ that depend on ratios of transport coefficients to
relaxation times.  In IS, causality is only established for the
\emph{linearized} theory around equilibrium.

The BDNK conditions~\eqref{eq:BDNK-eta-cond}--\eqref{eq:BDNK-coupled-ineq}
replace the IS conditions entirely.  The key differences are:
\begin{enumerate}
  \item No relaxation times appear; causality comes from the
    interplay of transport coefficients and frame coefficients.
  \item The causality proof is \emph{nonlinear}---it holds for
    arbitrary solutions, not only perturbations around equilibrium.
  \item The conditions are \emph{coupled}: one cannot check causality
    of the heat, bulk, and shear sectors independently.  This is
    physically natural, since in a relativistic fluid all dissipative
    channels are coupled through the energy-momentum tensor.
\end{enumerate}

\subsubsection{Conformal Fluid (Simplified Case)}
For a conformal fluid ($\zeta = 0$, $\varepsilon = 3p$,
$\cs^{2} = c^{2}/3$), the BDNK causality conditions simplify
considerably:
\begin{equation}
  \eta > 0\,,\qquad
  \text{frame coefficients chosen such that }
  v_{\text{char}}^{2} \le c^{2}\,,
  \label{eq:BDNK-conformal}
\end{equation}
where $v_{\text{char}}$ denotes the characteristic speeds of the
coupled sound--shear system.  This case is particularly relevant for
the quark-gluon plasma and ultrarelativistic astrophysical
environments.


%------------------------------------------------------------------
\subsection{Relativistic Buoyancy and Convective Stability}
\label{subsec:rel-buoyancy}

The classical theory of convective stability (Chandrasekhar, Chapters
II and IV) rests on the Schwarzschild and Ledoux criteria.  Their
relativistic generalizations involve the enthalpy density $w$ in place
of the mass density $\rho$.

\subsubsection{Relativistic Schwarzschild Criterion}
A fluid element displaced adiabatically in a gravitational field is
convectively unstable when the ambient entropy gradient is
superadiabatic.  In the relativistic formulation, the criterion
becomes
\begin{equation}
  \frac{\mathrm{d} s}{\mathrm{d} r}
  < \frac{s + p/T}{w}\,
    \frac{\mathrm{d} p}{\mathrm{d} r}\,
    \left(\frac{\partial s}{\partial p}\right)_{\!\varepsilon}
  \label{eq:rel-schwarzschild}
\end{equation}
(for a radially stratified equilibrium with $\mathrm{d} p/\mathrm{d} r < 0$).
The key relativistic modification is the appearance of $w = \varepsilon + p$
in the denominator, which enhances the effective inertia and thus
\emph{stabilizes} relativistically hot matter against convection
compared to the classical prediction.

\subsubsection{Relativistic Ledoux Criterion}
When the fluid has multiple species (e.g., baryon number and lepton
number in neutron-star matter), a composition gradient provides an
additional stabilizing influence.  The relativistic Ledoux criterion is
\begin{equation}
  \frac{\mathrm{d} s}{\mathrm{d} r}
  < \left(\frac{\mathrm{d} s}{\mathrm{d} r}\right)_{\!\text{ad}}
  + \frac{1}{w}\left(\frac{\partial p}{\partial Y_{e}}\right)_{\!\varepsilon,s}
    \frac{\mathrm{d} Y_{e}}{\mathrm{d} r}\,,
  \label{eq:rel-ledoux}
\end{equation}
where $Y_{e}$ is the electron fraction (or any composition variable)
and the subscript ``ad'' denotes the adiabatic gradient.

\subsubsection{Relativistic Brunt--V\"ais\"al\"a Frequency}
The squared Brunt--V\"ais\"al\"a frequency governing the oscillation
(or growth) of displaced fluid elements generalizes to
\begin{equation}
  N^{2}
  = -\frac{c^{2}}{w}\,\frac{\mathrm{d} p}{\mathrm{d} r}
    \left[
      \frac{1}{\cs^{2}}\,\frac{\mathrm{d} p}{\mathrm{d} r}
      - \frac{\mathrm{d}\varepsilon}{\mathrm{d} r}
    \right],
  \label{eq:rel-BV}
\end{equation}
where all quantities are evaluated in the equilibrium background.
Convective instability corresponds to $N^{2} < 0$.  In the
non-relativistic limit ($w \to \rho c^{2}$, $\varepsilon \to \rho c^{2}$),
Eq.~\eqref{eq:rel-BV} reduces to the standard expression
$N^{2} = -g\,(\mathrm{d}\ln\rho/\mathrm{d} r - (1/\gamma)(\mathrm{d}\ln p/\mathrm{d} r))$.

\subsubsection{Connection to Chandrasekhar's Thermal Instability}
Chandrasekhar's analysis of thermal instability (B\'enard convection)
proceeds from the Boussinesq equations.  The relativistic
generalization requires:
\begin{enumerate}
  \item Replacing $\rho$ by $w/c^{2}$ in the buoyancy term and in the
    inertia;
  \item Using the BDNK causal heat
    equation~\eqref{eq:BDNK-heat} in place of Fourier's law---note
    that unlike the Israel--Stewart telegraph equation, the BDNK
    treatment keeps the system first-order while still ensuring causal
    heat propagation;
  \item Including the $O(v^{2}/c^{2})$ corrections to the thermal
    expansion coefficient:
    $\alpha_{\text{rel}} = \alpha\,[1 + O(\varepsilon_{\text{th}}/\rho_{0}c^{2})]$;
  \item Verifying that the resulting critical Rayleigh number satisfies
    $\text{Ra}_{\text{crit,rel}} = \text{Ra}_{\text{crit,class}}
    \times [1 + O(v^{2}/c^{2})]$.
\end{enumerate}
The dominant relativistic effect on thermal instability is the
\emph{finite propagation speed of heat}, which in the BDNK framework
arises from the coupled characteristic speeds of the full
conservation system rather than from an independent thermal wave speed.
This is discussed in detail in the companion sections on relativistic
B\'enard convection.


%------------------------------------------------------------------
\subsection{Summary of Causality Bounds}
\label{subsec:causality-summary}

Table~\ref{tab:causality-bounds} collects all signal speeds and their
relativistic bounds in the BDNK framework.

\begin{table}[htbp]
\centering
\caption{Signal speeds in BDNK relativistic dissipative hydrodynamics.
  All characteristic speeds must satisfy $v \le c$.  The third column shows the
  modification relative to the classical (acausal) theory.  Unlike Israel--Stewart,
  the dissipative signal speeds are not given by simple formulas involving
  relaxation times but emerge from the characteristic polynomial of the full
  first-order system.}
\label{tab:causality-bounds}
\begin{tabular}{llll}
\hline\hline
\textbf{Mode} & \textbf{Signal speed (BDNK)} & \textbf{Classical limit} &
  \textbf{Causality bound} \\
\hline
Sound &
  $\cs = \sqrt{(\partial p/\partial\varepsilon)_{s}}$ &
  $\sqrt{\gamma p/\rho}$ &
  $\cs \le c$ \\[4pt]
Thermal &
  From characteristic polynomial &
  $\infty$ (Fourier) &
  Bounded by $c$ via BDNK \\[2pt]
  &
  (coupled to sound \& shear) &
  &
  frame conditions \\[4pt]
Bulk viscous &
  From characteristic polynomial &
  $\infty$ (Navier--Stokes) &
  Bounded by $c$ via BDNK \\[2pt]
  &
  (coupled to sound) &
  &
  frame conditions \\[4pt]
Shear viscous &
  From characteristic polynomial &
  $\infty$ (Navier--Stokes) &
  Bounded by $c$ via BDNK \\[2pt]
  &
  (decoupled in linearized theory) &
  &
  frame conditions \\[4pt]
Alfv\'en &
  $\vA = c\sqrt{b^{2}/(4\pi w + b^{2})}$ &
  $B/\sqrt{4\pi\rho}$ &
  $\vA \le c$ \\[4pt]
Fast magnetosonic &
  $v_{f}$ (see Eq.~52 of conventions) &
  $\sqrt{\cs^{2}+\vA^{2}}$ &
  $v_{f} \le c$ \\
\hline\hline
\end{tabular}
\end{table}

\subsubsection{Classical-to-Relativistic Modification Summary}
The key modifications that the BDNK relativistic framework imposes on
each classical transport equation are:
\begin{enumerate}
  \item \textbf{Fourier $\to$ BDNK (heat):}
    Replace Fourier's law with the BDNK constitutive
    relation~\eqref{eq:BDNK-heat} in a general hydrodynamic frame;
    causality follows from the coupled frame conditions, not from a
    relaxation term.
  \item \textbf{Navier--Stokes (bulk) $\to$ BDNK:}
    The constitutive relation $\Pi_{\text{bulk}} = -\zeta\,\nabla_{\mu}u^{\mu}$
    is retained at first order; causality is ensured by the BDNK frame
    coefficients acting on the full system.
  \item \textbf{Navier--Stokes (shear) $\to$ BDNK:}
    The constitutive relation $\Pi^{\mu\nu} = -2\eta\,\sigma^{\mu\nu}$
    is retained at first order; no relaxation time $\tau_{\pi}$ is needed.
  \item \textbf{Mass density $\to$ enthalpy density:}
    Replace $\rho$ by $w/c^{2} = (\varepsilon+p)/c^{2}$ in all inertial
    terms.
  \item \textbf{Sound speed:}
    Replace $\gamma p/\rho$ by $(\partial p/\partial\varepsilon)_{s}$;
    enforce $\le c^{2}$.
  \item \textbf{Buoyancy:}
    Replace Schwarzschild/Ledoux criteria with versions using $w$.
\end{enumerate}

\subsubsection{BDNK Causality Verification Checklist}
\label{subsubsec:causality-checklist}
For any relativistic stability calculation in the BDNK framework,
the following checks must be performed:
\begin{enumerate}
  \item[$\checkmark$] $\cs^{2} = (\partial p/\partial\varepsilon)_{s}$ satisfies
    $0 \le \cs^{2} \le c^{2}$.
  \item[$\checkmark$] Shear viscosity is strictly positive: $\eta > 0$.
  \item[$\checkmark$] Bulk--shear combination is positive:
    $\zeta + \tfrac{2}{3}\eta > 0$.
  \item[$\checkmark$] BDNK frame coefficients
    $(\varepsilon_{1}, \beta_{1}, \alpha_{1})$ satisfy the coupled
    strong-hyperbolicity inequalities~\eqref{eq:BDNK-coupled-ineq}.
  \item[$\checkmark$] All characteristic speeds from the principal symbol
    satisfy $|v_{\text{char}}| \le c$.
  \item[$\checkmark$] Entropy production rate
    $\nabla_{\mu}S^{\mu} \ge 0$
    (guaranteed by $\eta, \zeta, \kappa \ge 0$).
  \item[$\checkmark$] If MHD: Alfv\'en speed $\vA < c$; fast magnetosonic
    speed $v_{f} < c$.
  \item[$\checkmark$] Non-relativistic limits recover the classical
    Chandrasekhar equations (BDNK frame corrections vanish at
    $O(v^{2}/c^{2})$).
\end{enumerate}
\noindent
Failure of any of these checks indicates either an inadmissible EOS,
incorrect transport coefficients, or an invalid choice of BDNK frame
coefficients.


%------------------------------------------------------------------
\subsection{Application: Equation of State Comparison}
\label{subsec:eos-application}

Figure~\ref{fig:eos-comparison} compares the squared sound speed
$\cs^2/c^2$ as a function of energy density for three representative
equations of state relevant to relativistic stability analysis.

The \textbf{ideal relativistic gas} ($\Gamma$-law) interpolates between
the non-relativistic regime ($\cs^2 \ll c^2$) at low energy density
and the ultrarelativistic limit ($\cs^2 \to c^2/3$ for $\Gamma = 4/3$)
at high energy density.

The \textbf{QGP} (quark-gluon plasma) equation of state, inspired by
lattice QCD calculations \citep{BazavovEtAl2014,BorsanyiEtAl2014},
shows a characteristic dip in $\cs^2$ near the QCD crossover transition
at $\varepsilon \approx \varepsilon_{\mathrm{nuc}}$, before recovering
toward the conformal limit $\cs^2/c^2 = 1/3$ at high energy density.
This non-monotonic behaviour has significant consequences for the
stability of QGP flows in heavy-ion collisions and early-universe
dynamics.

The \textbf{nuclear matter} equation of state (APR-like stiff EOS
\citep{AkmalPR1998}) shows steady stiffening with increasing density,
with $\cs^2/c^2$ reaching $\sim 0.8$ at several times nuclear
saturation density.  Observational constraints from neutron star
masses \citep{LattimerPrakash2016,RileyEtAl2021} require such
stiffening to support stars above $2\,M_\odot$.  The causality
bound $\cs^2 \le c^2$ places a fundamental upper limit on the
stiffness of nuclear matter.

All three EOS satisfy the causality constraint $\cs^2/c^2 \le 1$
and the thermodynamic stability condition $\cs^2 \ge 0$ (Le
Chatelier's principle, eq.~\ref{eq:le-chatelier}).

\begin{figure}[htbp]
  \centering
  \includegraphics[width=0.85\textwidth]{fig_eos_comparison.pdf}
  \caption{Squared sound speed $\cs^2/c^2$ as a function of energy
    density (in units of nuclear saturation energy density
    $\varepsilon_{\mathrm{nuc}} \approx 150$\,MeV/fm$^3$) for three
    representative equations of state: ideal relativistic gas (blue),
    quark-gluon plasma (brown), and nuclear matter (pink).  The
    causality bound $\cs^2 = c^2$ (red line) and the conformal limit
    $\cs^2/c^2 = 1/3$ (gray dotted) are shown.  The QCD crossover
    produces a pronounced dip in the QGP sound speed near
    $\varepsilon \approx \varepsilon_{\mathrm{nuc}}$.}
  \label{fig:eos-comparison}
\end{figure}

%------------------------------------------------------------------
\subsection{Nuclear EOS Transport Coefficients for BDNK}
\label{subsec:nuclear-transport}

The practical application of BDNK hydrodynamics to neutron star
physics requires knowledge of the transport coefficients---shear
viscosity $\eta$, bulk viscosity $\zeta$, and thermal conductivity
$\kappa$---as functions of baryon density $n_B$ and temperature $T$.
These coefficients depend sensitively on the nuclear equation of state
through the composition (proton fraction $x_p$), effective nucleon
masses $m^*$, and the threshold for fast (direct Urca) neutrino
emission processes.

We compare transport coefficients computed for three representative
nuclear EOSs that span the range of current theoretical predictions:

\begin{itemize}
  \item \textbf{APR} \citep{AkmalPR1998}: A variational calculation
    with realistic two- and three-nucleon potentials.  Predicts a stiff
    EOS with $\cs^2/c^2$ reaching $\sim 0.8$ at high density.
  \item \textbf{SLy} \citep{DouchinHaensel2001}: A Skyrme-type
    effective interaction fitted to nuclear properties and neutron
    matter calculations.  Moderate stiffness.
  \item \textbf{DD2}: A relativistic mean-field model with
    density-dependent meson-nucleon couplings.  Softer at intermediate
    densities but stiffens at high density to support massive NSs.
\end{itemize}

The key features are:

\medskip\noindent
\textbf{Shear viscosity.}  In degenerate neutron star matter at
$T \sim 10^9$\,K, the shear viscosity is dominated by
neutron-neutron scattering and scales as
$\eta \sim n_B^{9/4}\, T^{-2}\, (m^*/m)^{-5/2}$
\citep{AlfordMahmoodi2012}.  The EOS enters primarily through the
effective mass $m^*(n_B)$, which decreases with density in all three
models but at different rates: DD2 predicts the strongest decrease,
giving the largest $\eta$ at high density.

\medskip\noindent
\textbf{Bulk viscosity.}  The bulk viscosity arises from the
$\beta$-equilibration lag during compression and expansion.  At low
density, modified Urca processes dominate, giving
$\zeta \propto n_B^2\, T^6$.  Above the direct Urca threshold
($x_p > x_{\mathrm{DU}} \approx 0.11$--$0.14$ depending on the EOS),
the bulk viscosity increases by several orders of magnitude
\citep{GavassinoAntonelli2024}.  DD2 reaches this threshold at lower
density than APR or SLy, producing a dramatic enhancement visible in
Figure~\ref{fig:nuclear-transport}.

\medskip\noindent
\textbf{Thermal conductivity.}  Electron thermal conduction dominates
in NS cores, with $\kappa \propto n_e^{1/3}\, T^{-1}$
\citep{PotekhinChabrier2018}.  The electron fraction $x_p$ determines
$n_e = x_p\, n_B$, so EOSs with higher proton fraction (DD2) predict
larger thermal conductivity.

\begin{figure}[htbp]
  \centering
  \includegraphics[width=\textwidth]{fig_nuclear_transport.pdf}
  \caption{BDNK transport coefficients as functions of baryon density
    $n_B/n_0$ (where $n_0 = 0.16$\,fm$^{-3}$ is nuclear saturation
    density) at $T = 10^9$\,K for three nuclear EOSs: APR (red solid),
    SLy (blue dashed), and DD2 (green dash-dotted).
    (a)~Shear viscosity $\eta$: dominated by nn-scattering, sensitive
    to the effective mass.
    (b)~Bulk viscosity $\zeta$: shows dramatic enhancement at the
    direct Urca threshold (highest density peak in DD2).
    (c)~Thermal conductivity $\kappa$: governed by electron conduction,
    sensitive to the proton fraction.
    The vertical dotted line marks nuclear saturation density.}
  \label{fig:nuclear-transport}
\end{figure}

These EOS-dependent transport coefficients enter directly into the
BDNK hyperbolicity conditions (Section~\ref{subsec:BDNK-causality}):
the frame coefficients $(\varepsilon_1, \beta_1, \alpha_1)$ must be
chosen in relation to $(\eta, \zeta, \kappa)$ and the thermodynamic
state $(\varepsilon, p, \cs^2)$ to ensure strong hyperbolicity.
The variation of transport coefficients across EOSs therefore affects
not only the magnitude of dissipative effects but also the
\emph{allowed range of BDNK frame choices} for a given physical
environment.

%------------------------------------------------------------------
\subsection{Transport Bounds from Stability and Causality}
\label{subsec:transport-bounds}

A deep connection between thermodynamic stability and causality
constraints on transport coefficients was established by Gavassino
\citep{GavassinoBounds2023}.  The key result is the fundamental
inequality
\begin{equation}
  \mathfrak{Im}\,\omega(k) \;\le\; |\mathfrak{Im}\,k|\,,
  \label{eq:fundamental-bound}
\end{equation}
which must hold for all dispersion relations $\omega(k)$ of a
covariantly stable system.  The proof is elegantly simple: if
\eqref{eq:fundamental-bound} were violated for some complex $k$, a
Lorentz boost with velocity
$v = \mathfrak{Im}\,k/\mathfrak{Im}\,\omega$ (subluminal since
$|v| < 1$ by assumption) would produce a frame in which a real-$k$
Fourier mode grows exponentially---an instability.  Thus covariant
stability \emph{alone} implies~\eqref{eq:fundamental-bound}.

This bound was originally derived by Heller, Serantes, Spali\'nski \&
Withers from microcausality (the vanishing of spacelike commutators
in quantum field theory).  Gavassino's proof shows that the same
bounds on transport emerge from requiring the fluid to be stable in
all inertial frames, without invoking any quantum-field-theoretic
assumptions.  The two derivations are complementary, reflecting the
deep connection between causality and covariant stability in
relativistic systems.

From \eqref{eq:fundamental-bound}, several concrete bounds follow
\citep{GavassinoBounds2023}:
\begin{enumerate}
  \item The infrared sound speed cannot exceed $c$:
    $c_{s} = \omega'(0) \le c$.
  \item The diffusion coefficient is non-negative and bounded above:
    $0 \le D \le 16/(3\pi K)$, where $K$ is the radius of convergence
    of the dispersion relation.
  \item Fluids with luminal sound speed ($c_{s} = c$) necessarily have
    \emph{vanishing} viscosities: $\eta = \zeta = 0$.
\end{enumerate}

\subsection{Dispersion Relations and Causality}
\label{subsec:dispersions-causality}

A related but distinct result by Gavassino, Disconzi \& Noronha
\citep{GavassinoDispersions2024} shows that individual dispersion
relations cannot, by themselves, guarantee causality.  Any single
dispersion branch $\omega(k)$ (other than the trivial $\omega = a + bk$)
necessarily produces superluminal tails when used alone to construct wave
packets.  This is because compactly-supported initial data have
entire-function Fourier transforms, but $\omega(k)$ satisfying
\eqref{eq:fundamental-bound} is \emph{never} an entire function (unless
linear in $k$).

The resolution is that causality is a \emph{collective} property
of all excitation branches: when the full initial-value problem is
solved---with all branches superposed---the superluminal tails from
different branches cancel exactly, provided the theory is local and
covariantly stable.  This cancellation is analogous to the need for
antiparticles in relativistic quantum mechanics to ensure causal
propagation.

For the stability analyses in this book, this result has an
important implication: the non-hydrodynamic modes present in
Israel--Stewart theory (or in any second-order formulation) are not
mere artifacts but play an essential physical role in cancelling
superluminal tails from the hydrodynamic modes, thereby ensuring
causal signal propagation.  In the BDNK framework, the same
cancellation occurs through the coupling of hydrodynamic modes
mediated by the frame coefficients, without introducing extra
dynamical degrees of freedom.

%------------------------------------------------------------------
\subsection{Causality Constraints on Radiative Transfer}
\label{subsec:causal-rad-transfer}

A striking application of the transport bounds discussed above arises
in radiation hydrodynamics.  \citet{GavassinoCausalRadTransfer2025}
showed that Spiegel's classical formula for the smoothing of
temperature fluctuations by radiative transfer---widely used in
stellar structure calculations---is \emph{unstable under Lorentz
boosts}.  The instability arises because Spiegel's derivation neglects
the finite transit time of photons, allowing the transport coefficients
to fall outside the ``hydrohedron'' (the convex set of
causality-compatible transport coefficients identified by Heller,
Serantes, Spali\'nski \& Withers).

Gavassino resolves this by solving the linearized radiative transfer
equations \emph{exactly} (accounting for the finite speed of photons),
and demonstrates that:
\begin{enumerate}
  \item The exact dispersion relation for heat waves takes the form
    $\omega(k) = -i\,(3D_h/\tau^2)[1 - \arctan(k\tau)/(k\tau)]$,
    which naturally interpolates between diffusive ($k\tau \ll 1$) and
    free-streaming ($k\tau \gg 1$) behavior;
  \item All infinite transport coefficients arising from the exact
    solution fall \emph{inside} the hydrohedron, as required by
    relativistic causality;
  \item The linearized radiation transfer equations are causal and
    covariantly stable \emph{by construction}, since the Boltzmann
    equation with emission and absorption respects the speed-of-light
    limit at the kinetic level.
\end{enumerate}
This result is directly relevant to the convective stability analyses
of Chapters~II and~VI, where radiation-mediated thermal diffusion
sets the effective thermal conductivity entering the Rayleigh number.

%------------------------------------------------------------------
%% =====================================================================
%%  Gavassino integration: Thermoelectric conduction in GR
%%  Paper 45: arXiv:2509.23845
%%  Paper 44: arXiv:2508.08936 — Charge separation in accelerating plasmas
%%  Paper 48: arXiv:1906.03140 — Thermodynamics of uncharged multifluids
%% =====================================================================
\subsection{Relativistic Thermoelectricity and Gravitoelectric Effects}
\label{subsec:rel-thermoelectricity}

In non-relativistic solid-state physics, the coupling between heat
and charge transport is described by the Seebeck and Peltier effects
within the standard Onsager framework.  The relativistic
generalisation of thermoelectricity is non-trivial because (i) heat
carries inertia ($E = mc^2$ implies that a heat flux is also a
momentum density), and (ii) gravitational redshift modifies the
equilibrium conditions for both temperature and chemical potential.

\begin{relcorrection}
\textbf{Causal, stable relativistic thermoelectricity.}
Gavassino~\cite{GavassinoThermoelectric2025} has constructed the first
fully causal, stable, and well-posed relativistic theory of
thermoelectric transport in rigid solids embedded in curved
spacetime, employing the BDNK strategy.  The key results are:
\begin{enumerate}
  \item \emph{Ohm--Fourier--Seebeck constitutive relations:}
    The electric current $J^\mu$ and heat flux $q^\mu$ in a Born-rigid
    conductor with Killing vector $\mathcal{K}^\mu = \mathcal{K}\,u^\mu$
    satisfy coupled transport equations with electrical conductivity
    $\sigma$, Seebeck coefficient $\mathcal{S}$, Peltier coefficient
    $\Pi = T\mathcal{S}$ (Onsager reciprocity), and thermal conductivity
    $\kappa$.  In the Lorenz gauge, the resulting initial-value problem
    is causal and locally well-posed in the fully nonlinear regime.

  \item \emph{Gravitothermoelectric effects:}
    Three novel relativistic phenomena emerge:
    (a)~acceleration-driven charge separation (Schiff--Barnhill and
    Stewart--Tolman effects);
    (b)~non-uniform Joule heating due to time dilation across an
    accelerating circuit;
    (c)~gravitationally redshifted magnetic diffusion.

  \item \emph{Relativistic Thomas--Fermi equation:}
    For compact objects, the equilibrium charge distribution satisfies a
    relativistic Thomas--Fermi equation that accounts for both the
    Tolman--Ehrenfest temperature gradient and Seebeck-driven charge
    displacements from cooling.
\end{enumerate}
These results demonstrate that the BDNK formalism extends naturally
beyond simple heat conduction to the full thermoelectric problem,
maintaining the causal and stable properties that are the hallmark
of the framework presented in this chapter.
\end{relcorrection}

\begin{relcorrection}
\textbf{Charge separation in accelerating relativistic plasmas.}
Gavassino~\cite{GavassinoChargeSeparation2025} has derived from first
principles the equilibrium electrothermal stratification of a fully
ionised plasma under irrotational relativistic acceleration in curved
spacetime.  The analysis reveals that for strong accelerations
(acceleration number
$\mathscr{A} = L_{\mathrm{Debye}}/L_{\mathrm{accel}} \gtrsim 1$),
significant charge separation develops, and the equilibrium electric
field is \emph{independent of the conductivity} of the medium.  Near
black hole horizons, $\mathscr{A}$ can approach unity, invalidating
the standard MHD charge-neutrality assumption.  This has direct
implications for the stability analyses in Chapters~X--XI, where MHD
equilibria serve as background states.
\end{relcorrection}

\begin{relcorrection}
\textbf{Thermodynamics of multifluid mixtures.}
The thermodynamic framework underlying these transport phenomena was
established by Gavassino \&
Antonelli~\cite{GavassinoAntonelliMultifluidThermo2019}, who
developed a gauge-invariant thermodynamic formalism for relativistic
multifluids with arbitrary numbers of superfluid components.  The
key insight is that the equation of state in thermodynamic
equilibrium contains fewer degrees of freedom than needed for the
Lagrangian (variational) approach, revealing a gauge freedom that
simplifies the hydrodynamic formulation.  This formalism provides
the thermodynamic underpinning for both the thermoelectric transport
theory and the multifluid stability analyses in later chapters.
\end{relcorrection}

\subsection*{References for This Section}

\begin{itemize}
  \item Bemfica, Disconzi \& Noronha, \textit{Phys.\ Rev.\ D}
    \textbf{98}, 104064 (2018) \citep{BemficaDN2018}
  \item Kovtun, \textit{JHEP} \textbf{10}, 034 (2019) \citep{Kovtun2019}
  \item Hoult \& Kovtun, \textit{JHEP} \textbf{06}, 067 (2020)
    \citep{HoultKovtun2020}
  \item Bemfica, Disconzi, Noronha \& Kovtun, \textit{Phys.\ Rev.\ D}
    \textbf{107}, 076012 (2023) \citep{BemficaDNK2023}
  \item Israel \& Stewart, \textit{Ann.\ Phys.}\ \textbf{118}, 341
    (1979) \citep{IsraelStewart1979}
  \item Rezzolla \& Zanotti, \textit{Relativistic Hydrodynamics},
    Oxford (2013) \citep{RezzollaZanotti2013}
  \item Akmal, Pandharipande \& Ravenhall, \textit{Phys.\ Rev.\ C}
    \textbf{58}, 1804 (1998) \citep{AkmalPR1998}
  \item Lattimer \& Prakash, \textit{Phys.\ Rep.}\ \textbf{621}, 127
    (2016) \citep{LattimerPrakash2016}
  \item Riley et~al., \textit{Astrophys.\ J.\ Lett.}\ \textbf{918},
    L27 (2021) \citep{RileyEtAl2021}
  \item Bazavov et~al., \textit{Phys.\ Rev.\ D} \textbf{90}, 094503
    (2014) \citep{BazavovEtAl2014}
  \item Bors\'anyi et~al., \textit{Phys.\ Lett.\ B} \textbf{730}, 99
    (2014) \citep{BorsanyiEtAl2014}
  \item Weinberg, \textit{Gravitation and Cosmology}, Wiley (1972)
    \citep{Weinberg1972}
  \item Chandrasekhar, \textit{Hydrodynamic and Hydromagnetic Stability},
    Oxford (1961) \citep{Chandrasekhar1961}
  \item Pandya \& Pretorius, \textit{Phys.\ Rev.\ D} \textbf{104},
    023015 (2021) \citep{PandyaPretorius2021}
  \item Pandya, Most \& Pretorius, arXiv:2209.09265 (2022)
    \citep{PandyaMostPretorius2022ideal}
  \item Gavassino \& Antonelli, \textit{Phys.\ Lett.\ B} \textbf{849},
    138434 (2024) \citep{GavassinoAntonelli2024}
  \item Hoult \& Kovtun, \textit{Phys.\ Rev.\ D} \textbf{106},
    066023 (2022) \citep{HoultKovtun2022kinetic}
  \item Potekhin \& Chabrier, \textit{Astron.\ Astrophys.}\ \textbf{609},
    A74 (2018) \citep{PotekhinChabrier2018}
  \item arXiv:2509.15303, NS evolution with BDNK (2025)
    \citep{NSevolutionBDNK2025}
  \item Gavassino, \textit{Phys.\ Lett.\ B} \textbf{840}, 137854
    (2023) \citep{GavassinoBounds2023}
  \item Gavassino, Disconzi \& Noronha, \textit{Phys.\ Rev.\ Lett.}\
    \textbf{132}, 162301 (2024) \citep{GavassinoDispersions2024}
  \item Gavassino, arXiv:2509.23845, thermoelectric conduction in GR
    (2025) \citep{GavassinoThermoelectric2025}
  \item Gavassino, arXiv:2508.08936, charge separation in accelerating
    plasmas (2025) \citep{GavassinoChargeSeparation2025}
  \item Gavassino \& Antonelli, \textit{MNRAS} \textbf{494},
    3882 (2020) \citep{GavassinoAntonelliMultifluidThermo2019}
\end{itemize}
